% chapters/tikz_ch10.tex - 第 10 章全部插图的高精度原生 TikZ 矢量代码

% 1. 图 10.1: 单个纤维内的径向集 R_pm 与 S*R^n 中的零双特征线
\newcommand{\drawFigTenOne}{%
\begin{tikzpicture}[scale=0.92, >=Stealth, line join=round]
    % 左图: T*_{(0,y)}R^n 的单个纤维
    \begin{scope}
        \draw[->] (-2.75, 0) -- (2.75, 0) node[right] {$\xi$};
        \draw[->] (0, -1.95) -- (0, 1.95) node[above] {$\eta$};
        % 径向集 R_+
        \draw[very thick] (0.22, 0) -- (2.40, 0);
        % 径向集 R_-
        \draw[very thick] (-2.40, 0) -- (-0.22, 0);
        \node[above] at (1.55, 0.10) {$\mathcal{R}_+$};
        \node[above] at (-1.55, 0.10) {$\mathcal{R}_-$};
        % H_p = -d/dxi 的流 (箭头指向 xi 减小的方向)
        \draw[->, thick, gray!75] (2.10, -0.32) -- (1.60, -0.32);
        \draw[->, thick, gray!75] (-1.60, -0.32) -- (-2.10, -0.32);
        \fill (0, 0) circle (1.6pt) node[above left, inner sep=1pt] {$o$};
        \node[below] at (0, -0.62) {\footnotesize $H_p = -\partial_\xi$};
    \end{scope}
    % 右图: S*_{(0,y)}R^n 中的零双特征线
    \begin{scope}[xshift=7.3cm]
        \draw[thick] (0, 0) circle (1.78);
        \fill (1.78, 0) circle (2.4pt);
        \fill (-1.78, 0) circle (2.4pt);
        % 流: 沿圆周趋向左侧的边界点
        \draw[->, thick] (0, 1.78) arc (90:137:1.78);
        \draw[->, thick] (0, -1.78) arc (270:223:1.78);
        \node[right] at (1.92, 0) {$\partial\mathcal{R}_+$};
        \node[left] at (-1.92, 0) {$\partial\mathcal{R}_-$};
        \node[below] at (0, -1.95) {\footnotesize $S^*_{(0,y)}\R^n$};
    \end{scope}
\end{tikzpicture}%
}

% 2. 图 10.2: 径向集 R 作为 H_p 的法向源
%    Char(P) 在 R 处与 S*M 相切 (R 为 Char(P) cap S*M 的孤立点),
%    H_p 沿 Char(P) 与沿纤维径向都从 R 出发 (法向源)
\newcommand{\drawFigTenTwo}{%
\begin{tikzpicture}[scale=1.0, >=Stealth, line join=round]
    % 纤维径向内部
    \fill[gray!8] (-3.15, 0) rectangle (3.15, -2.10);
    % 余球面丛 (紧化余切丛的边界)
    \draw[very thick] (-3.15, 0) -- (3.15, 0);
    \node[right] at (3.15, 0) {$S^*M$};
    \node at (-2.55, -1.76) {$\overline{T}^*M$};
    % 特征集 Char(P): 与 S*M 切于 R 的弧 (远段)
    \draw[thick] (1.55, -0.19) .. controls (2.20, -0.35) .. (2.75, -0.62);
    \draw[thick] (-1.55, -0.19) .. controls (-2.20, -0.35) .. (-2.75, -0.62);
    \node[right] at (2.78, -0.62) {$\Char(P)$};
    % H_p 沿 Char(P) 自 R 流出的箭头 (即 Char(P) 的近段)
    \draw[->, very thick] (0.05, -0.01) .. controls (0.55, -0.10) and (1.05, -0.14) .. (1.52, -0.18);
    \draw[->, very thick] (-0.05, -0.01) .. controls (-0.55, -0.10) and (-1.05, -0.14) .. (-1.52, -0.18);
    % H_p 沿纤维径向自 R 流出
    \draw[->, very thick] (0.06, -0.26) -- (0.36, -1.68);
    \node[right] at (0.42, -1.66) {$\tilde{H}_p$};
    % 径向集 (单点)
    \fill (0, 0) circle (2.8pt) node[above left, inner sep=3pt] {$\mathcal{R}$};
\end{tikzpicture}%
}

% 3. 图 10.3: 定理 10.7 的图示 (流出 R / 流入 R)
\newcommand{\drawFigTenThree}{%
\begin{tikzpicture}[scale=0.88, >=Stealth, line join=round]
    % 左图: 数据在 WF'(B), 后向流趋向 R
    \begin{scope}
        % Char(P) cap S*M 附近的特征带
        \fill[gray!15] (0, -0.92) -- (3.05, -0.92) -- (3.05, 0.02) -- (0, 0.02) -- cycle;
        \draw[thin] (0, -0.92) -- (3.05, -0.92);
        \draw[very thick] (-3.10, 0.02) -- (3.10, 0.02);
        \node[above] at (-2.55, 0.06) {\footnotesize $S^*M$};
        % Ell(G)
        \draw[dashed] (0.10, -1.00) rectangle (1.35, 0.10);
        \node[below] at (0.72, -1.04) {\footnotesize $\Ell(G)$};
        % WF'(B)
        \draw[dashed] (2.10, -1.00) rectangle (3.02, 0.10);
        \node[above] at (2.56, 0.14) {\footnotesize $\WF'(B)$};
        % 后向流: 指向 R
        \draw[->, very thick] (2.02, -0.45) -- (1.42, -0.45);
        \draw[->, very thick] (1.05, -0.45) -- (0.32, -0.45);
        % 径向集
        \fill (0, 0.02) circle (2.6pt) node[above left, inner sep=2pt] {$\mathcal{R}$};
    \end{scope}
    % 右图: 前向流自 R 出发进入 Ell(E)
    \begin{scope}[xshift=7.60cm]
        \fill[gray!15] (0, -0.92) -- (3.05, -0.92) -- (3.05, 0.02) -- (0, 0.02) -- cycle;
        \draw[thin] (0, -0.92) -- (3.05, -0.92);
        \draw[very thick] (-3.10, 0.02) -- (3.10, 0.02);
        \node[above] at (-2.55, 0.06) {\footnotesize $S^*M$};
        \draw[dashed] (0.10, -1.00) rectangle (1.35, 0.10);
        \node[below] at (0.72, -1.04) {\footnotesize $\Ell(B) \supset \WF'(B)$};
        \draw[dashed] (2.10, -1.00) rectangle (3.02, 0.10);
        \node[above] at (2.56, 0.14) {\footnotesize $\Ell(E)$};
        % 前向流: 离开 R
        \draw[->, very thick] (0.32, -0.45) -- (1.05, -0.45);
        \draw[->, very thick] (1.42, -0.45) -- (2.02, -0.45);
        \fill (0, 0.02) circle (2.6pt) node[above left, inner sep=2pt] {$\mathcal{R}$};
    \end{scope}
\end{tikzpicture}%
}
